\documentclass[dvipdfmx]{jsarticle}
\usepackage{graphicx}
\usepackage{url}

\title{ソフトウェア設計及び実験\\第6回（UI2）}
\author{学生番号 \ 6125010806　氏名 \ 清水 \ 統貴}
\date{2026年5月27日提出}

\begin{document}
\maketitle

\section{課題の目的}
%
本課題の目的は以下のとおりである．\\
・ゲームを彩るUIのプログラミングを学び, ゲーム開発に必要な基本的なサウンド処理, 
アニメーション処理 (, 当たり判定)の知識・スキルを身につける.\\

\section{原理}
%
\subsection{アニメーション原理}
%
アニメーションの原理としてはパラパラの要領で, 複数枚の画像を用意して, 一定の
時間間隔で次々に表示するというものである. 今回のプログラム内では\verb+SDL_AddTimer+
が画像を次々に送る役割をしている. 一定間隔ごとに画面を消し, 次のコマを作り直すという
処理を繰り返すことで実現している. キャラクターの描画ではキャラクターをスプライトから描画
しており, 画像の必要な部分だけを切り取って描画している. 
\section{課題1}
%

\section{感想}
%


\begin{thebibliography}{99}
%
% \bibitem{chatGPT} Open AI : ChatGPT \url{https://chatgpt.com/ja-JP} (2026年4月30日閲覧).
% %
% \bibitem{claude} Anthropic : Claude \url{https://claude.ai/new} (2026年5月5日閲覧).
% %
\bibitem{gemini} Google : Gemini \url{https://gemini.google.com/?hl=ja}
%
\bibitem{qiita} IronNyan : 画像を２値化する【道具としての OpenCV 画像認識】 \url{https://qiita.com/IronNyan/items/644feed2f16925e0ee06}
%
\bibitem{qiita} awesomecity : OpenCV　「画像のグレースケール化」 \url{https://qiita.com/awesomecity1998/items/448a3f91a489c8959d20}
%
\bibitem{nsi} NSI : 2値化とは？ \url{https://nsi-freak.com/binarization/#:~:text=2%E5%80%A4%E5%8C%96%E3%81%A8%E3%81%AF%E3%80%81%E7%94%BB%E5%83%8F%E3%82%92%E7%99%BD%E3%81%A8%E9%BB%92,%E3%80%8C%E9%96%BE%E5%80%A4%E3%80%8D%E3%81%A8%E8%A8%80%E3%81%84%E3%81%BE%E3%81%99%E3%80%82&text=%E3%83%A2%E3%83%8E%E3%82%AF%E3%83%AD%E7%94%BB%E5%83%8F%E3%81%AF0%EF%BD%9E,%E3%81%A8%E5%AE%9A%E7%BE%A9%E3%81%95%E3%82%8C%E3%81%A6%E3%81%84%E3%81%BE%E3%81%99%E3%80%82}
%
\bibitem{kuroro} 黒光俊秀 : OpenCVで使われるcannyとは?利用法からCanny法の理論などを徹底解説 \url{https://kuroro.blog/python/wOt3yEohr7oQt1qzif71/}
%
\end{thebibliography}
\end{document}